\documentclass[10pt,aspectratio=169]{beamer}

% ---------- Кодировка и язык ----------
\usepackage[utf8]{inputenc}
\usepackage[T2A]{fontenc}
\usepackage[english,russian]{babel}

% ---------- Оформление ----------
\usetheme{default}
\usecolortheme{seahorse}
\useinnertheme{rectangles}
\setbeamertemplate{navigation symbols}{}
\setbeamertemplate{footline}[frame number]
\setbeamertemplate{theorems}[numbered]
\setbeamertemplate{blocks}[rounded=false]
\setbeamerfont{block title}{series=\bfseries}
\setbeamertemplate{itemize items}[circle]
\setbeamertemplate{enumerate items}[default]

% ---------- Математика ----------
\usepackage{amsmath,amssymb,mathtools}
\usepackage{bm}

% ---------- Окружения с русскими названиями и отдельными счётчиками ----------
\uselanguage{Russian}
\languagepath{Russian}
\deftranslation[to=Russian]{Proof}{Доказательство}
\theoremstyle{plain}
\newtheorem{thm}{Теорема}
\newtheorem{lem}{Лемма}
\newtheorem{prop}{Утверждение}
\newtheorem{cor}{Следствие}
\theoremstyle{definition}
\newtheorem{defn}{Определение}

% ---------- Макросы (как в notes.tex) ----------
\newcommand{\R}{\mathbb{R}}
\newcommand{\N}{\mathbb{N}}
\newcommand{\abs}[1]{\left\lvert #1 \right\rvert}
\newcommand{\set}[1]{\left\{ #1 \right\}}
\newcommand{\ip}[2]{\langle #1, #2\rangle}
\newcommand{\Rpp}[1][2]{\R^{#1}_{++}}
\newcommand{\Rp}[1][2]{\R^{#1}_{+}}
\DeclareMathOperator{\conv}{conv}
\DeclareMathOperator{\cone}{cone}
\DeclareMathOperator{\interior}{int}
\DeclareMathOperator{\Sp}{Sp}
\newcommand{\PhiHat}{\hat{\Phi}}

\title{Неоклассическая разрешимость вектора выпусков\\ при ценах в общем положении}
\author{Олег Лаврин}
\institute{}
\date{}

\begin{document}

\begin{frame}
    \titlepage
\end{frame}

% =====================================================================
\section{Постановка задачи}

\begin{frame}{Неоклассические функции и разрешимость}
    \begin{defn}
        Функция $h:\Rp[n]\to\R_{+}$ называется \emph{неоклассической}, если она непрерывна, вогнута и однородна степени~1. Будем обозначать класс неоклассических функций n переменных как $\Phi_n$.
        Неоклассическая функция $h$ называется \emph{строго положительной}, если 
        $\forall p \ne 0$ значение $h(p)>0$.
        Класс строго положительных функций $n$ переменных обозначим $\PhiHat_n$.
    \end{defn}


    \begin{defn}[неоклассическая разрешимость]
        Пусть $\set{\hat p_t\in\Rpp}_{t\in[T]}$~--- некоторый конечный набор цен. Будем говорить, что вектор выпусков $y\in\Rpp[T]$ \emph{неоклассически разрешим} при этих ценах, если существует такая неоклассическая функция $h\in\Phi_2$ и такие окрестности $\set{U_t\ni\hat p_t}_{t\in[T]}$, что для каждого набора $\set{p_t\in U_t}_{t\in[T]}$ существует полная абсолютно непрерывная мера $\mu$ на $\Rp$, что
        \[
            \forall t\in[T]\Rightarrow\mu(A_t)=y_t,\qquad\text{где } A_t=\set{x\in\Rpp\mid h(p_t\circ x)<1}.
        \]
    \end{defn}
\end{frame}


% =====================================================================
\section{Общее положение}


\begin{frame}{Трансверсальность}
    \begin{defn}[трансверсальное пересечение]\label{def:transversal}
        Пусть $h,g\in\PhiHat_2$, $x_0\in\Rpp: h(x_0)=g(x_0)=1$. Будем говорить, что линии уровня $h(x)=1$ и $g(x)=1$ пересекаются в $x_0$ \emph{трансверсально}, если $\partial h(x_0)\cap\partial g(x_0)=\varnothing$.
    \end{defn}

    \begin{defn}[общее положение относительно $h$]\label{def:gp}
        Пусть $h\in\PhiHat_2$ и $P=\set{p_t\in\Rpp}_{t\in[T]}$. Пара $(h,P)$ находится \emph{в общем положении}, если
        \begin{enumerate}
            \item для любой пары различных индесов $s\neq t$ линии уровня $h(p_{s}\circ x)=1$ и $h(p_{t}\circ x)=1$ имеют конечное кол-во пересечений, причем все эти пересечения трансверсальны.
            \item ни для какой тройки различных индексов $r\neq s\neq t\neq r$ линии уровня
            \\ $h(p_{r}\circ x)=1, h(p_{s}\circ x)=1$ и $h(p_{t}\circ x)=1$ не пересекаются в одной точке.
        \end{enumerate}
    \end{defn}

    \emph{Замечание 1.} Это определение равносильно трансверсальности кривых $h(x)=1$ и $g(x)=1$ в точке $x_0$ в смысле негладкого анализа (Kruger, 2006; Ioffe, 2017).

    \emph{Замечание 2.} Для $h(x)=\sqrt{x^1x^2}\exp\bigl(\tfrac14\sin\ln\tfrac{x^1}{x^2}\bigr)$ пересечение кривых $h(x)=1$ и $h(p\circ x)=1$ счётно на открытом множестве $p$, потому нужна строгая \mbox{положительность}.
\end{frame}

\begin{frame}{Общее положение пары $(h,P)$}
    \begin{lem}\label{lem:finite}
        Пусть $h\in\PhiHat_2$ и $G_h\subset\Rpp$~--- множество всех цен $p$, при которых кривые $h(x)=1$ и $h(p\circ x)=1$ имеют конечное число точек пересечения, и все они трансверсальны.\\
        Тогда $G_h$ открыто и всюду плотно в $\Rpp$.
    \end{lem}

    \begin{prop}[устойчивость общего положения]
        Пусть $h\in\PhiHat_2$, $P\in\Rpp[2\times T]$ и пара $(h,P)$ находится в общем положении. Тогда существует такая окрестность $V\subset\Rpp[2\times T]$ набора $P$, что для каждого набора цен $Q\in V$ пара $(h,Q)$ также находится в общем положении.
    \end{prop}

    \begin{prop}[плотность общего положения]
        Пусть $h\in\PhiHat_2$. Тогда в любом непустом открытом множестве $U\subset\Rpp[2\times T]$ найдётся набор цен $P$, для которого пара $(h,P)$ находится в общем положении.
    \end{prop}

    \begin{block}{Замечание}
        Утверждения~\ref{prop:stable} и~\ref{prop:dense} показывают корректность определения~\ref{def:gp}. Общее положение не нарушается при малом изменении цен и достигается сколь угодно малым изменением.
    \end{block}
\end{frame}

\setcounter{prop}{0}
\begin{frame}{Устойчивость общего положения}
    \begin{prop}[Устойчивость общего положения]\label{prop:stable}
        Пусть $h\in\hat\Phi_{2}, P\in\Rpp[2\times T]$ и пара $(h,P)$ находится в общем положении.
        Тогда существует такая окрестность $V\subset\Rpp[2\times T]$ набора $P$, что для каждого набора
        цен $Q\in V$ пара $(h,Q)$ также находится в общем положении.
    \end{prop}

    \begin{proof}
        Пункт~1 определения~\ref{def:gp} сохраняется вблизи $P$ по лемме~\ref{lem:finite}: для пары $(s,t)$ он означает $p_t/p_s\in G_h$, а $G_h$ открыто.
        Пусть пункт~2 нарушается для наборов $Q_n\to P$, каждый раз для одних и тех же индексов $r,s,t$. Введём множество $I_{st}(Q)$ --- точек пересечения линий уровня $h(q_s\circ x)=1$ и $h(q_t\circ x)=1$ в $\Rp$ для набора цен $Q=\{q_t\}_{t\in[T]}$.

        При ценах, близких к $P$, все точки пересечения лежат в одном ограниченном множестве: из $1=h(q_s\circ x)\ge h(1,0)q_s^1x^1+h(0,1)q_s^2x^2$ следует $x^1\le1/(h(1,0)q_s^1)$ и $x^2\le1/(h(0,1)q_s^2)$, а $q_s^1,q_s^2$ отделены от нуля. Поэтому из тройных точек $x_n\in I_{rs}(Q_n)\cap I_{rt}(Q_n)$ можно выбрать сходящуюся подпоследовательность, и её предел лежит в $I_{rs}(P)\cap I_{rt}(P)$ по непрерывности $h$, то есть является тройной точкой для $P$, что противоречит общему положению $(h,P)$.
    \end{proof}
\end{frame}

\begin{frame}{Плотность общего положения}
    \begin{prop}[Плотность общего положения]\label{prop:dense}
        Пусть $h\in\hat{\Phi}_{2}$. Тогда в любом непустом открытом множестве $U\subset\Rpp[2\times T]$
        найдётся набор цен $P$, для которого пара $(h,P)$ находится в общем положении.
    \end{prop}

    \begin{proof}
        Для пары $(s,t)$ пункт~1 определения~\ref{def:gp} означает $p_t/p_s\in G_h$. Обозначим через $W_{st}$ множество наборов $P\in\Rpp[2\times T]$, для которых он выполнен. Отображение $P\mapsto p_t/p_s$ непрерывно, поэтому $W_{st}$ открыто как прообраз открытого по лемме~\ref{lem:finite} множества $G_h$. Множество $W_{st}$ плотно: возьмём любой набор $P$ и любую его окрестность. По лемме~\ref{lem:finite} найдётся $q\in G_h$, сколь угодно близкое к $p_t/p_s$. Заменив $p_t$ на $p_s\circ q$, получим набор из этой окрестности, лежащий в $W_{st}$. Конечное пересечение открытых плотных множеств открыто и плотно, поэтому $V:=U\bigcap_{s\ne t}W_{st}$ --- непустое открытое множество.
    \end{proof}
\end{frame}

\begin{frame}{Плотность общего положения: окончание доказательства}
    \begin{proof}[Доказательство (окончание)]
        Возьмём $P\in V$ и пройдясь по индексам $k\in[T]$ будем по очереди заменять $p_k$ на $\lambda_kp_k$ с $\lambda_k$ близким к $1$, не выходя из открытого множества $V$. По однородности кривая $k$ при этом растягивается из начала координат в $1/\lambda_k$ раз, а остальные кривые не меняются. Покажем по индукции, что после $k$-го шага пункт~2 выполнен для всех троек с индексами не большими $k$. Тройки с индексами меньше $k$ шаг $k$ не затрагивает. Для тройки $(r,s,k)$, $r,s<k$, множество $I_{rs}$ конечно по пункту~1 и от $\lambda_k$ не зависит, а кривая $k$ проходит через точку $x\in I_{rs}$ лишь при $\lambda_k=1/h(p_k\circ x)$. Значит, плохих $\lambda_k$ конечное число, и мы всегда можем выбрать подходящее $\lambda_k$, и после $T$-го шага получаем набор $P\in V\subset U$ что пара $(h,P)$ находится в общем положении.
    \end{proof}
\end{frame}

\begin{frame}{Общее положение цен}
    \begin{defn}[Общее положение цен]\label{def:prices-gp}
        Будем говорить, что набор цен $P:=\set{p_t\in\Rpp}_{t\in[T]}$ находится в
        \textit{общем положении}, если при $s\neq t$ верно, что $p^{1}_{s}\neq p^{1}_{t}$ и $p^{2}_{s}\neq p^{2}_{t}$.
    \end{defn}

    \begin{prop}\label{prop:prices-gp}
        Набор цен $P$ находится в общем положении тогда и только тогда, когда существует $h\in\hat\Phi_{2}$, для которой пара $(h,P)$ находится в общем положении.
    \end{prop}

    \begin{proof}
        \emph{Необходимость.} Пусть пара $(h,P)$ в общем положении. Если $p^{1}_{s}=p^{1}_{t}$ при $s\ne t$, то точка $(1/(p^{1}_{s}h(1,0)),0)$ лежит на кривых $h(p_s\circ x)=1$ и $h(p_t\circ x)=1$ и не лежит в $\Rpp$, а значит, не является трансверсальной, что противоречит пункту~1 определения~\ref{def:gp}; так же для вторых координат.
    \end{proof}
\end{frame}

\begin{frame}{Общее положение цен: доказательство достаточности}
    \begin{proof}[Доказательство (окончание)]
        \emph{Достаточность.} Зафиксируем непрерывную строго выпуклую функцию $f\colon[0,a]\to[0,b]$ с $f(0)=b>0$ и $f(a)=0$. Зафиксируем квадратную решётку с шагом $\epsilon>0$ и в каждой ячейке, пересекающейся с графиком $f$, выберем по точке графика. Соединим соседние точки отрезками, получив ломаную, заданную функцией $g(x):=\min_j\ip{\xi_j}{x}\in\hat\Phi_2$, где $\ip{\xi_j}{x}=1$ --- прямые, содержащие её звенья. Поскольку цены в общем положении, общие точки кривых $s\ne t$ не лежат на осях: общая точка на оси означала бы $p^1_s=p^1_t$ или $p^2_s=p^2_t$. По той же причине при достаточно малом $\epsilon$ они лежат на разных звеньях. Если при этом никакие три прямые $\ip{\xi_i\circ p_r}{x}=1$, $\ip{\xi_j\circ p_s}{x}=1$, $\ip{\xi_k\circ p_t}{x}=1$ с различными $i,j,k$ не проходят через одну точку и никакие две не совпадают, то пересечения кривых трансверсальны, конечны и происходят только по парам. Оба условия выражаются одним неравенством $\det\bigl(\xi_i\circ p_r-\xi_j\circ p_s,\ \xi_i\circ p_r-\xi_k\circ p_t\bigr)\neq0,$
        так как при совпадении двух прямых первый столбец нулевой. Таких ограничений конечное число, и каждое --- ненулевой многочлен от $\xi$: при $\xi_i=(1,0)$, $\xi_j=(0,1)$, $\xi_k=0$ определитель равен $p^1_rp^2_s\neq0$. Поэтому в любой окрестности набора $\xi$ есть набор, где все они выполнены, и сколь угодно малым изменением $\xi_j$, сохраняющим вершины ломаной внутри их ячеек, пара $(g,P)$ приводится в общее положение.
    \end{proof}
\end{frame}

% =====================================================================
\section{Регулярная разрешимость}

\begin{frame}{Регулярная разрешимость и спектр}
    \begin{defn}[регулярная разрешимость]
        Пусть $P:=\set{p_t\in\Rpp}_{t\in[T]}$~--- набор цен. Будем говорить, что вектор выпусков $y\in\Rpp[T]$ \emph{регулярно разрешим} при ценах $P$, если существуют функция $h\in\PhiHat_2$, для которой пара $(h,P)$ находится в общем положении, и полная абсолютно непрерывная мера $\mu$ на $\Rp$, для которой $\mu(A_t)=y_t$ при всех $t\in[T]$, где $A_t:=\set{x\in\Rpp\mid h(p_t\circ x)<1}$.
    \end{defn}

    \begin{defn}[спектр]\label{def:spectrum}
        Пусть $h\in\Phi_2$ и $P:=\set{p_t\in\Rpp}_{t\in[T]}$~--- набор цен. Каждой точке $x\in\Rpp$ сопоставим множество индексов $S(x):=\set{t\in[T]\mid h(p_t\circ x)<1}$ и для $S\subset[T]$ положим $R_S:=\set{x\in\Rpp\mid S(x)=S}$. \emph{Спектром} пары $(h,P)$ назовём семейство
        \[
            \Sp(h,P):=\set{S\subset[T]\mid\interior R_S\neq\varnothing}.
        \]
    \end{defn}

    Множества $R_S$ --- слои отображения $x\mapsto S(x)$, поэтому $\Rpp=\bigsqcup\limits_{S\subset[T]}R_S$.
\end{frame}

\begin{frame}{Мера и конус спектра}
    \begin{lem}\label{lem:cone}
        Рассмотрим неоклассическую функцию $h\in\Phi_2$ и набор цен $P:=\set{p_t\in\Rpp}_{t\in[T]}$. Тогда вектор выпусков $y\in\Rpp[T]$ принадлежит конусу $\cone\set{\mathbf 1_S\mid S\in\Sp(h,P)}$, если и только если существует полная абсолютно непрерывная мера $\mu$ на $\Rp$, для которой при всех $t\in[T]$ выполнено интегральное условие $\mu(A_t)=y_t$, где $A_t:=\set{x\in\Rpp\mid h(p_t\circ x)<1}$.
    \end{lem}

    \begin{proof}
        Множества $R_S$ разбивают $\Rpp$, и $A_t=\bigsqcup_{S\ni t}R_S$, поэтому условия $\mu(A_t)=y_t$ означают $y=\sum\mu(R_S)\mathbf 1_S$. При $S\notin\Sp(h,P)$ множество $R_S$ без точек кривых $h(p_t\circ x)=1$ открыто и потому пусто, а сами кривые имеют нулевую меру как границы выпуклых множеств, поэтому $\mu(R_S)=0$. Обратно, разложение $y=\sum m_S\mathbf 1_S$ задаёт массу $m_S\ge0$, которую можно абсолютно непрерывно распределить по произвольному шару в $\interior R_S$.
    \end{proof}
\end{frame}

\begin{frame}{Приближение функцией в общем положении}
    \begin{lem}\label{lem:approx}
        Пусть набор цен $P$ находится в общем положении и $h\in\Phi_2$. Тогда существует
        $g\in\hat\Phi_2$, для которой пара $(g,P)$ находится в общем положении и
        $\mathrm{Sp}(g,P)\supseteq\mathrm{Sp}(h,P)$.
    \end{lem}

    \begin{proof}
        Выберем для каждого $S\in\mathrm{Sp}(h,P)$ точку $x_S\in\operatorname{int} R_S$, и введем множество $Z:=\set{p_t\circ x_S \mid t\in[T],S\in\mathrm{Sp}(h,P)}$ получив \textit{зазор} $\eta:=\min_{z\in Z}\abs{h(z)-1}>0$. Функция $H(x):=h(x)+\varepsilon\bigl(x^1+x^2+\tfrac{x^1x^2}{x^1+x^2}\bigr)$ при достаточно малом $\varepsilon>0$ отличается от $h$ на $Z$ меньше чем на $\eta/2$, а значит для каждого $z\in Z$ верно что $H(z)\gtrless1\Leftrightarrow h(z)\gtrless1$. Более того, $H\in\hat\Phi_2$ и строго вогнута на отрезке $x^1+x^2=1$, так что её линия уровня $H(x)=1$ --- график строго выпуклой убывающей функции. Выберем на этой линии точки от одной оси до другой, соединим соседние отрезками и положим $g(x):=\min_j\ip{\xi_j}{x}$, где $\ip{\xi_j}{x}=1$ --- прямые, содержащие эти отрезки. Тогда $g\in\hat\Phi_2$, линия $g(x)=1$ --- полученная ломаная, и если точки расположены достаточно часто, то $\abs{g-H}<\eta/2$ на $Z$. Сдвигая векторы $\xi_m$ на сколь угодно малые величины, приведём пару $(g,P)$ в общее положение, сохранив это неравенство. Тогда $\abs{g-h}<\eta$ на $Z$, поэтому $g(p_t\circ x_S)-1$ и $h(p_t\circ x_S)-1$ одного знака, и по непрерывности $g$ точка $x_S$ имеет спектр $S$ а значит $S\in\mathrm{Sp}(g,P)$.
    \end{proof}
\end{frame}

\begin{frame}{Теорема 1: неоклассическая и регулярная разрешимость}
    \begin{thm}\label{thm:regular}
        Пусть набор цен $P$ находится в общем положении. Тогда вектор выпусков $y\in\Rpp[T]$
        неоклассически разрешим при ценах $P$ тогда и только тогда, когда он регулярно разрешим
        при этих ценах.
    \end{thm}

    \begin{proof}
        По лемме~\ref{lem:cone} в обоих случаях речь идёт о включении $y$ в конус, натянутый на булевы векторы спектра. Пусть $y$ регулярно разрешим с функцией $h\in\hat\Phi_2\subset\Phi_2$. Во внутренности каждой области $R_S$, $S\in\mathrm{Sp}(h,P)$, есть точка вне кривых, а строгие неравенства в ней сохраняются при малом изменении цен; поэтому $\mathrm{Sp}(h,Q)\supseteq\mathrm{Sp}(h,P)$ для всех $Q$ из некоторой окрестности $P$, и $y$ неоклассически разрешим. Обратно, если $y$ неоклассически разрешим, возьмём $h\in\Phi_2$ и меру при самих ценах $P$; функция $g$ из леммы~\ref{lem:approx} находится в общем положении с $P$ и $\mathrm{Sp}(g,P)\supseteq\mathrm{Sp}(h,P)$, так что $y$ регулярно разрешим.
    \end{proof}
\end{frame}

% =====================================================================
\section{Достижимые спектры}

\begin{frame}{Достижимые спектры}
    \begin{defn}[достижимые спектры]
        Пусть $P:=\set{p_t\in\Rpp}_{t\in[T]}$~--- набор цен. Будем говорить, что $t$ \emph{покрывает} пару $s,r$, если $p_t\in\conv\set{p_s,p_r}+\Rp$. Будем говорить, что $S\subset[T]$~--- \emph{достижимый спектр} при ценах $P$, если ни один индекс $t\in S$ не покрывает пару индексов $s,r\notin S$. Будем обозначать как $\Sp(P)$ семейство таких спектров.
    \end{defn}

    \begin{lem}\label{lem:reach-nec}
        Пусть $h\in\Phi_{2}$ и $P:=\{p_{t}\in\Rpp\}_{t\in[T]}$~--- набор цен. Тогда $\mathrm{Sp}(h,P)\subseteq\mathrm{Sp}(P)$.
    \end{lem}

    \begin{proof}
        Предположим противное: пусть $S\in\mathrm{Sp}(h,P)$ и $t\in S$ покрывает пару $s,r\notin S$. Но тогда $A_{t}\subset A_{s}\cup A_{r}$, а значит $s$ или $r$ лежит в $S$, что противоречит условию.
    \end{proof}
\end{frame}

\begin{frame}{Реализация одного спектра}
    \begin{lem}\label{lem:one}
        Пусть $P:=\{p_{t}\in\Rpp\}_{t\in[T]}$~--- набор цен. Тогда для любого $S\in\mathrm{Sp}(P)$ существует кусочно-линейная $h(x):=\min\limits_{k\in K}\langle\xi_{k}, x\rangle\in\hat\Phi_{2}$ для которой выполнено что $S\in\mathrm{Sp}(h,P)$.
    \end{lem}

    \begin{proof}
        Зафиксируем $S\in\mathrm{Sp}(P)$ и построим для каждого $t\in S$ вектор $\xi_t\in\Rpp$ что $\ip{\xi_t}{p_t}<1$ и $\ip{\xi_t}{p_s}>1$ при всех $s\notin S$. Тогда $h(x):=\min\limits_{t\in S}\ip{\xi_t}{x}$ лежит в $\hat\Phi_2$, $h(p_t)\le\ip{\xi_t}{p_t}<1$ при $t\in S$ и $h(p_s)=\min\limits_{t\in S}\ip{\xi_t}{p_s}>1$ при $s\notin S$. Неравенства строгие, поэтому по непрерывности точка $x=(1,1)$ лежит в $\interior R_{S}$ со спектром $S\in\mathrm{Sp}(h,P)$.

        Приведём пример таких $\xi_t$. Зафиксируем $t\in S$. Для $s\notin S$ индекс $t$ не покрывает пару $s,s$, то есть $p_s\not\le p_t$: выполнено $p^1_s>p^1_t$ или $p^2_s>p^2_t$. Пусть $\sigma_t>0$ таково, что $\sigma_t\,(p^1_s-p^1_t)+(p^2_s-p^2_t)>0$ для всех $s\notin S$. Тогда $\ip{(\sigma_t,1)}{p_t}<\min_{s\notin S}\ip{(\sigma_t,1)}{p_s}$, и вектор
        \[
            \xi_t:=\frac{2\,(\sigma_t,1)}{\ip{(\sigma_t,1)}{p_t}+\min_{s\notin S}\ip{(\sigma_t,1)}{p_s}}
        \]
        удовлетворяет $\ip{\xi_t}{p_t}<1<\ip{\xi_t}{p_s}$ при всех $s\notin S$.
    \end{proof}
\end{frame}

\begin{frame}{Реализация одного спектра: окончание доказательства}
    \begin{proof}[Доказательство (окончание)]
        Убедимся, что такое $\sigma_t$ существует. Неравенство $\sigma_t\,(p^1_s-p^1_t)+(p^2_s-p^2_t)>0$ при $p^1_s=p^1_t$ выполнено всегда (тогда $p^2_s>p^2_t$), при $p^1_s>p^1_t$ означает $\sigma_t>\frac{p^2_t-p^2_s}{p^1_s-p^1_t}$, а при $p^1_s<p^1_t$~--- $\sigma_t<\frac{p^2_s-p^2_t}{p^1_t-p^1_s}$. Поэтому достаточно проверить, что для любых $r,s\notin S$ с $p^1_r>p^1_t>p^1_s$
        \[
            \frac{p^2_t-p^2_r}{p^1_r-p^1_t}<\frac{p^2_s-p^2_t}{p^1_t-p^1_s},
        \]
        и взять $\sigma_t$ больше нуля и всех левых частей, но меньше всех правых. Точка $q:=\lambda p_s+(1-\lambda)p_r$ с $\lambda:=\frac{p^1_r-p^1_t}{p^1_r-p^1_s}\in(0,1)$ лежит на отрезке $[p_s,p_r]$ и $q^1=p^1_t$; так как $t$ не покрывает пару $s,r$, имеем $q\not\le p_t$, то есть $q^2>p^2_t$. Подставляя $q^2=\lambda p^2_s+(1-\lambda)p^2_r$, получаем $\lambda(p^2_s-p^2_t)>(1-\lambda)(p^2_t-p^2_r)$, то есть $(p^1_r-p^1_t)(p^2_s-p^2_t)>(p^1_t-p^1_s)(p^2_t-p^2_r)$, и после деления на $(p^1_r-p^1_t)(p^1_t-p^1_s)>0$ это и есть проверяемое неравенство.
    \end{proof}
\end{frame}

\begin{frame}{Склейка спектров}
    \begin{lem}[склейка]\label{lem:patch}
        Пусть $P:=\{p_{t}\in\Rpp\}_{t\in[T]}$~--- набор цен и $S_1,\dots,S_k\in\mathrm{Sp}(P)$. Тогда существует кусочно-линейная $h\in\hat\Phi_2$, для которой $S_1,\dots,S_k\in\mathrm{Sp}(h,P)$.
    \end{lem}

    \begin{proof}
        Поскольку $\varnothing\in\mathrm{Sp}(h,P)$ для любой $h\in\hat\Phi_2$, пустые $S_i$ можно отбросить; считаем все $S_i$ непустыми. Для каждой пары $(i,t)$ с $t\in S_i$ возьмём по доказательству предыдущей леммы вектор $\xi_{i,t}\in\Rpp$ с $\ip{\xi_{i,t}}{p_t}<1$ и $\ip{\xi_{i,t}}{p_s}>1$ при $s\notin S_i$. Выберем $M>0$ так, чтобы $M\min(\xi^1_{i,t}p^1_s,\ \xi^2_{i,t}p^2_s)>1$ для всех таких пар $(i,t)$ и всех $s\in[T]$, и числа $0<\alpha_1<\dots<\alpha_k$ с $\alpha_{i+1}\ge M\alpha_i$. Положим $x_i:=(\alpha_i,\alpha_i^{-1})$ и $\tilde\xi_{i,t}:=(\xi^1_{i,t}\alpha_i^{-1},\ \xi^2_{i,t}\alpha_i)$. Тогда $\ip{\tilde\xi_{i,t}}{p\circ x_i}=\ip{\xi_{i,t}}{p}$ для любой цены $p$, а при $l\ne i$
        \[
            \ip{\tilde\xi_{i,t}}{p\circ x_l}=\xi^1_{i,t}p^1\,\frac{\alpha_l}{\alpha_i}+\xi^2_{i,t}p^2\,\frac{\alpha_i}{\alpha_l}\ \ge\ M\min(\xi^1_{i,t}p^1,\ \xi^2_{i,t}p^2)>1,
        \]
        поскольку одно из отношений $\alpha_l/\alpha_i$, $\alpha_i/\alpha_l$ не меньше $M$, а оба слагаемых положительны.
    \end{proof}
\end{frame}

\begin{frame}{Склейка спектров: окончание доказательства}
    \begin{proof}[Доказательство (окончание)]
        Положим
        \[
            h(x):=\min_{i\in[k],\,t\in S_i}\ip{\tilde\xi_{i,t}}{x}\in\hat\Phi_2.
        \]
        Для $t\in S_i$ имеем $h(p_t\circ x_i)\le\ip{\tilde\xi_{i,t}}{p_t\circ x_i}=\ip{\xi_{i,t}}{p_t}<1$. Для $s\notin S_i$ каждый из векторов даёт в точке $p_s\circ x_i$ значение больше единицы: векторы $\tilde\xi_{i,t}$ той же копии~--- по выбору $\xi_{i,t}$, векторы других копий~--- по выбору $M$; значит $h(p_s\circ x_i)>1$. Все неравенства строгие, поэтому точка $x_i$ принадлежит камере $\interior R_{S_i}$ со спектром $S_i\in\mathrm{Sp}(h,P)$.
    \end{proof}
\end{frame}

\begin{frame}{Теорема 2: критерий разрешимости}
    \begin{thm}\label{thm:admissible-cone}
        Пусть набор цен $P$ находится в общем положении. Тогда вектор выпусков $y\in\Rpp[T]$ регулярно (равносильно, неоклассически) разрешим при ценах $P$ тогда и только тогда, когда $y\in\cone\set{\mathbf 1_S\mid S\in\Sp(P)}$.
    \end{thm}

    \begin{proof}
        Равносильность регулярной и неоклассической разрешимости~--- теорема~\ref{thm:regular}. Пусть $y$ регулярно разрешим с функцией $h\in\PhiHat_2$. По лемме~\ref{lem:cone} $y\in\cone\set{\mathbf 1_S\mid S\in\Sp(h,P)}$, а $\Sp(h,P)\subseteq\Sp(P)$ по лемме о достижимых спектрах. Обратно, пусть $y=\sum_{i\in[k]}m_i\mathbf 1_{S_i}$ с $m_i>0$ и $S_i\in\Sp(P)$. По лемме~\ref{lem:patch} найдётся $h\in\PhiHat_2$ с $S_1,\dots,S_k\in\Sp(h,P)$, по лемме~\ref{lem:approx}~--- функция $g\in\PhiHat_2$ в общем положении с $P$ и $\Sp(g,P)\supseteq\Sp(h,P)\ni S_1,\dots,S_k$, и по лемме~\ref{lem:cone} существует полная абсолютно непрерывная мера $\mu$ с $\mu(A_t)=y_t$ для $g$. Значит $y$ регулярно разрешим.
    \end{proof}
\end{frame}

\begin{frame}{Замечание}
    \begin{block}{Замечание}
        Для $P$ в общем положении существует $h\in\PhiHat_2$, для которой пара $(h,P)$ находится в общем положении и $\Sp(h,P)=\Sp(P)$. Действительно, применим лемму~\ref{lem:patch} ко всему семейству $\Sp(P)$ и получим $h_0\in\PhiHat_2$ с $\Sp(h_0,P)\supseteq\Sp(P)$; по лемме~\ref{lem:approx} найдётся $h\in\PhiHat_2$ в общем положении с $P$ и $\Sp(h,P)\supseteq\Sp(h_0,P)$, а обратное включение $\Sp(h,P)\subseteq\Sp(P)$ даёт лемма о достижимых спектрах. Тем самым каждый регулярно разрешимый вектор выпусков реализуется одной и той же функцией $h$, и от $y$ зависит только мера.
    \end{block}
\end{frame}

\begin{frame}{Алгоритм проверки}
    Теорема~\ref{thm:admissible-cone} даёт конечную процедуру проверки регулярной разрешимости $y$ при ценах $P$ в общем положении.
    \begin{enumerate}
        \item \emph{Построение $\mathrm{Sp}(P)$.} Для каждой тройки индексов $t,s,r$ (допускается $s=r$) проверяем, покрывает ли $t$ пару $s,r$: условие $p_t\in\operatorname{conv}\set{p_s,p_r}+\Rp$ означает $p_t\ge\lambda p_s+(1-\lambda)p_r$ при некотором $\lambda\in[0,1]$, то есть непустоту пересечения двух отрезков значений $\lambda$, по одному на каждую координату. Затем перебираем $S\subset[T]$ и оставляем те, в которых ни один $t\in S$ не покрывает пары $s,r\notin S$. Это $O(T^3)$ проверок покрытий и $2^T$ кандидатов.
        \item \emph{Проверка принадлежности конусу.} Решаем задачу линейного программирования: найти $m_S\ge0$, $S\in\mathrm{Sp}(P)$, с $\sum_Sm_S\mathbf 1_S=y$. Если она разрешима, базисное решение даёт разложение $y$ не более чем по $T$ спектрам; вместе с леммой~\ref{lem:patch} и леммой~\ref{lem:approx} это явная функция $h$ и мера, реализующие $y$. Если нет, двойственная задача даёт $\lambda\in\R^T$ с $\sum_{t\in S}\lambda_t\ge0$ при всех $S\in\mathrm{Sp}(P)$ и $\ip{\lambda}{y}<0$ (лемма Фаркаша); это опровержение: для любой реализующей пары $(h,\mu)$ из $y=\sum_S\mu(R_S)\mathbf 1_S$ и $\mathrm{Sp}(h,P)\subseteq\mathrm{Sp}(P)$ следовало бы $\ip{\lambda}{y}\ge0$. Неравенство $\sum_{\lambda_t<0}\abs{\lambda_t}y_t\le\sum_{\lambda_t>0}\lambda_ty_t$ выполнено для всех разрешимых $y$ и нарушено для данного.
    \end{enumerate}
\end{frame}

\begin{frame}{Заключение}
    Для набора цен $P$ в общем положении (попарно различные первые и вторые координаты):
    \begin{itemize}
        \item неоклассическая разрешимость (с окрестностями цен и произвольной $h\in\Phi_2$) равносильна регулярной разрешимости (одна строго положительная $h$, пара $(h,P)$ в общем положении);
        \item семейство спектров $\Sp(h,P)$ по всем $h\in\Phi_2$ описывается комбинаторно: $S$ достижим тогда и только тогда, когда ни один $t\in S$ не покрывает пары $s,r\notin S$;
        \item все достижимые спектры реализуются одной кусочно-линейной функцией в общем положении с $P$;
        \item $y$ разрешим тогда и только тогда, когда $y\in\cone\set{\mathbf 1_S\mid S\in\Sp(P)}$; проверка сводится к задаче линейного программирования с сертификатом Фаркаша в случае неразрешимости.
    \end{itemize}
\end{frame}

\end{document}
